% sec-06: the quantitative case.  Figures 16 (d2), 17 (graphviz).
\section{The Clinical Evidence a Funder Is Buying}

\subsection{Six Quantities a Reviewer Can Check}

Every quantity below is published or deposited, every one carries the limitation
its own authors stated, and no in-silico result is presented as a clinical one.
The discipline matters more here than anywhere else in the document, because a
capitalization plan that overstates its evidence is asking for money under a
false description.

\begin{apptable}
\begin{tabularx}{\textwidth}{>{\raggedright\arraybackslash}p{2.7cm} >{\raggedright\arraybackslash}p{1.9cm} >{\raggedright\arraybackslash}p{2.1cm} >{\raggedright\arraybackslash}p{1.4cm} >{\raggedright\arraybackslash}X}
\headrow \thead{Source} & \thead{Test arm} & \thead{Comparator} & \thead{Tier} & \thead{Limitation, in its own authors' words}\\
\midrule
RASolute 302, May 2026 & 13.2 mo mOS & 6.6 mo mOS & Trial & Metastatic and previously treated; silent on the resectable setting \\
QSP simulation, 10 arms, 250 ODEs & 12.8 mo mOS, HR 0.25 & 5.4 mo mOS & In silico & Assumes no acquired resistance and ideal pharmacodynamics; grade 3 plus high in all arms \\
Digital twin, 1000 patients & 12.1 mo mOS & not applicable & Twin & No patient-specific PK, PD or tumour growth parameters; simple Emax model \\
Digital twin, PFS hazard ratio & HR 0.31 & not applicable & Twin & Same source, same limitation; no immune compartments \\
VVUQ credibility, 55 tests & score 81.9 & V and V 40 gate & Twin & A pre-trial credibility score; not a post-trial validation \\
Empirical triplicate, 100,000 records & G3 plus 8.0 pct & G3 plus 25.0 pct & In silico & The G12C log favours experimental while the KRAS-mutant report favours control \\
\end{tabularx}
\end{apptable}
\vspace{-0.65cm}
\tabcap{Six quantities with their comparator, their evidence tier, and the
limitation each set of authors stated. No row can be read without the caveat
that qualifies it, which is the condition on citing any of them.}

\begin{appfloat}[!tb]
\begin{appfig}[d2-type / grid with interval measures]
% A seven by five grid on the left and an interval panel on the right,
% separated by a 10mm corridor that no element spans.  All four value columns
% take the identical 17mm width.
\foreach \c/\x/\lab in {0/1.65/{Quantity}, 1/3.90/{Test arm}, 2/5.70/{Comparator},
  3/7.50/{Tier}, 4/9.30/{Date}}{
  \ifnum\c=0
    \node[d2cellh,text width=31mm,minimum height=7.0mm] at (\x,0) {\lab};
  \else
    \node[d2cellh,minimum width=17mm,text width=14mm,minimum height=7.0mm] at (\x,0) {\lab};
  \fi}
\foreach \y/\q in {-0.76/{RASolute 302 median OS, RAS G12},
  -1.52/{QSP simulation median OS, 10 arms},
  -2.28/{Digital twin median OS, 1000 patients},
  -3.04/{Digital twin PFS hazard ratio},
  -3.80/{VVUQ credibility score, 55 tests},
  -4.56/{Empirical triplicate, grade 3 plus}}{
  \node[d2celll,text width=31mm,minimum height=7.0mm] at (1.65,\y) {\q};}
\foreach \y/\asty/\a/\b/\csty/\c/\d in {%
  -0.76/d2key/{13.2 months}/{6.6 months}/d2cell/{Trial}/{2026-05},
  -1.52/d2cellk/{12.8 months}/{5.4 months}/d2cellg/{In silico}/{2025-08},
  -2.28/d2cellk/{12.1 months}/{not applicable}/d2cellg/{Twin}/{2025-10},
  -3.04/d2cellk/{HR 0.31}/{not applicable}/d2cellg/{Twin}/{2025-10},
  -3.80/d2cellk/{score 81.9}/{V and V 40 gate}/d2cellg/{Twin}/{2025-10},
  -4.56/d2cellk/{8.0 percent}/{25.0 percent}/d2cellg/{In silico}/{2025-07}}{
  \node[\asty,minimum width=17mm,text width=14mm,minimum height=7.0mm] at (3.90,\y) {\a};
  \node[d2cell,minimum width=17mm,text width=14mm,minimum height=7.0mm] at (5.70,\y) {\b};
  \node[\csty,minimum width=17mm,text width=14mm,minimum height=7.0mm] at (7.50,\y) {\c};
  \node[d2cell,minimum width=17mm,text width=14mm,minimum height=7.0mm] at (9.30,\y) {\d};}
\begin{scope}[shift={(11.10,-4.95)}]
  \axisx{0}{4.4}{0}
  \foreach \m/\lab in {0/0, 1.25/5, 2.50/10, 3.75/15}{
    \draw[protogray,line width=0.3pt] (\m,0) -- (\m,-0.10);
    \node[font=\tiny\sffamily,text=protogray,anchor=north] at (\m,-0.12) {\lab};}
  \node[font=\tiny\sffamily,text=protogray,anchor=north] at (2.0,-0.42) {months, median overall survival};
  \ciband{3.65}{2.55}{4.05}{3.30}{13.2 test}
  \ciband{3.05}{1.15}{2.15}{1.65}{6.6 control}
  \legkey{0}{2.35}{pablue1}{12.8 QSP, point estimate}
  \legkey{0}{1.95}{pablue1}{12.1 twin, point estimate}
  \legkey{0}{1.55}{pagraym}{no interval: deterministic run}
  \ptitle{0}{4.45}{Intervals where a source reports one}
\end{scope}
\ptitle{0}{0.62}{Six checkable quantities}
\node[d2cellk,text width=24mm,minimum height=9mm] at (13.10,-5.85)
  {2.4-fold simulated in 2025 against 2.0-fold observed in 2026};
\node[font=\tiny\itshape,text=protogray,anchor=north west,text width=132mm]
  at (0,-6.55) {Three quantities carry an interval and three do not. A ten-arm
  deterministic ODE run and a 1000-patient synthetic cohort have no sampling
  interval, and drawing one would invent precision the source does not claim.};
\end{appfig}
\vspace{-0.65cm}
\figcaption{Six quantities a reviewer can check against a published source,\\
each beside its comparator and its stated limitation. Two are in\\
silico, one is a digital twin, and three are trial or registry.}
\end{appfloat}

In August 2025 a ten-arm quantitative systems pharmacology simulation returned
median overall survival of 12.8 months against 5.4 for chemotherapy
\cite{qspmetpancre}. In May 2026, nine months later, RASolute 302 reported 13.2
months against 6.6 in the RAS G12 population of previously treated metastatic
pancreatic cancer \cite{rasolute302}. The simulated ratio was 2.4-fold and the
observed ratio was 2.0-fold.

That proximity is a chronology observation and a hypothesis-supporting one. It
is \textbf{not} a validation claim, and three differences are material: 1000
simulated against 241 enrolled; a combination against a single agent; and KRAS
G12C against a primarily G12D and G12V population. Those clauses appeared in all
ten prior applications in the same paragraph as the numbers, and they appear
here for the same reason \cite{tenapps2026}.

\subsection{The Trial Being Funded}

Pancreatic ductal adenocarcinoma remains the disease with the widest gap between
incidence and survival in the United States \cite{siegel2025statistics}, and the
standard of care after resection is adjuvant combination chemotherapy
\cite{conroy2018folfirinox}. The candidate agent was selected from a June 2025
meta-analysis across forty literature areas \cite{daraxident}, and the
1000-patient digital twin that produced the hazard ratio of 0.31 and the
credibility score of 81.9 is deposited with its verification notebooks
\cite{daraxsims}.

The trial is open-label, single-arm, 3+3 escalation, up to 18 treated
participants, with perioperative daraxonrasib and a staged eight-arm robotic
pancreaticoduodenectomy in which an on-premises large language model is confined
to advisory output \cite{onpremwhippl,phase1ind}. The device specification a
reviewer will ask for is fixed in the protocol: 56 degrees of freedom, 640
sensor channels, a 3 N per-arm and 18 N cumulative force cap, a 3 ms cross-arm
arm stop against a 500 ms system stop, and a five-vessel no-fly gate. Dose
levels are 160, 220 and 300 mg with a 28-day dose-limiting toxicity window.

Daraxonrasib is investigational and is already in Phase 3 evaluation. It is
nowhere described in this paper as first in human, and the supportable novelty
claim concerns the integrated surgical and advisory workflow rather than the
agent. That correction was introduced during the audit pass on the ten prior
applications and is carried forward here unchanged.

\subsection{Where Each Number Lands}

\begin{appfloat}[!tb]
\begin{appfig}[graphviz-type / record chain]
% Four ranks left to right, one cluster per rank.  Exactly two targets receive
% two edges; each convergent pair is one straight edge and one bent at 18, so
% the arrivals are 4mm apart on the target's west face.
\foreach \i/\y/\a/\b/\c/\sty in {%
  1/0/{RASolute 302}/{13.2 vs 6.6 months}/{NEJM 2026}/gvkey,
  2/-1.55/{QSP simulation}/{12.8 vs 5.4 months}/{250 ODEs}/gvsoft,
  3/-3.10/{Digital twin}/{HR 0.31, score 81.9}/{55 tests}/gvsoft,
  4/-4.65/{Interlock bench}/{arm 3 ms, system 500 ms}/{200 runs}/gvsoft,
  5/-6.20/{Empirical triplicate}/{8.0 vs 25.0 percent}/{100,000 records}/gvsoft}{
  \node[\sty,text width=25mm] (q\i) at (0,\y) {\textbf{\a}\\\b\\\c};}
\begin{scope}[on background layer]
\node[gvcluster,fit=(q1)(q2)(q3)(q4)(q5),inner sep=6pt] (cq) {};
\end{scope}
\node[gvctitle,anchor=south west] at ([yshift=1mm]cq.north west) {Published quantity};
\foreach \i/\y/\a/\b in {7/-0.75/{IND 7}/{Previous human experience},
  6/-2.30/{IND 6}/{Pharmacology and toxicology},
  8/-3.85/{IND 8}/{Additional information},
  5/-5.40/{IND 5}/{Proposed clinical research}}{
  \node[gvboxs,text width=23mm] (n\i) at (4.35,\y) {\textbf{\a}\\\b};}
\begin{scope}[on background layer]
\node[gvcluster,fit=(n5)(n6)(n7)(n8),inner sep=6pt] (cn) {};
\end{scope}
\node[gvctitle,anchor=south west] at ([yshift=1mm]cn.north west) {IND section};
\foreach \i/\y/\a/\b in {1/0/{Protocol 1}/{Background},
  3/-1.55/{Protocol 3}/{Dose rationale},
  9/-3.10/{Protocol 9}/{Statistical considerations},
  6/-4.65/{Protocol 6}/{Device and interlocks},
  7/-6.20/{Protocol 7}/{Safety monitoring}}{
  \node[gvboxs,text width=23mm] (r\i) at (8.55,\y) {\textbf{\a}\\\b};}
\begin{scope}[on background layer]
\node[gvcluster,fit=(r1)(r3)(r6)(r7)(r9),inner sep=6pt] (cr) {};
\end{scope}
\node[gvctitle,anchor=south west] at ([yshift=1mm]cr.north west) {Protocol section};
\foreach \i/\y/\a/\b/\c in {5/-0.40/{M5}/{IND amendment}/{FDA acknowledgment},
  4/-2.30/{M4}/{VVUQ freeze}/{SHA-256 manifest},
  3/-4.20/{M3}/{Bench verification}/{p95 latency table},
  8/-6.10/{M8}/{Dose level 1 cleared}/{DSMB minutes}}{
  \node[gvboxg,text width=25mm] (m\i) at (12.75,\y) {\textbf{\a} \b\\\c};}
\begin{scope}[on background layer]
\node[gvcluster2,fit=(m3)(m4)(m5)(m8),inner sep=6pt] (cm) {};
\end{scope}
\node[gvctitle,text=protogray,anchor=south west] at ([yshift=1mm]cm.north west)
  {Milestone and artifact};
\draw[gvedgeb] (q1) -- node[font=\tiny,fill=protowhite,inner sep=1.2pt] {cited} (n7);
\draw[gvedge] (q2) -- node[font=\tiny,fill=protowhite,inner sep=1.2pt] {informs} (n6);
\draw[gvedge] (q3) -- node[font=\tiny,fill=protowhite,inner sep=1.2pt] {supports} (n8);
\draw[gvedge] (q4) -- node[font=\tiny,fill=protowhite,inner sep=1.2pt] {specifies} (n5);
\draw[gvedge] (q5.east) to[bend right=18]
  node[font=\tiny,fill=protowhite,inner sep=1.2pt,pos=0.6] {bounds} (n7.south west);
\draw[gvedge] (n7) -- node[font=\tiny,fill=protowhite,inner sep=1.2pt] {frames} (r1);
\draw[gvedge] (n6) -- node[font=\tiny,fill=protowhite,inner sep=1.2pt] {sets} (r3);
\draw[gvedge] (n8) -- node[font=\tiny,fill=protowhite,inner sep=1.2pt] {bounds} (r9);
\draw[gvedge] (n5) -- node[font=\tiny,fill=protowhite,inner sep=1.2pt] {fixes} (r6);
\draw[gvedge] (n7.south east) to[bend right=16]
  node[font=\tiny,fill=protowhite,inner sep=1.2pt,pos=0.65] {triggers} (r7.west);
\draw[gvedgeb] (r1) -- node[font=\tiny,fill=protowhite,inner sep=1.2pt] {filed in} (m5);
\draw[gvedge] (r3) -- node[font=\tiny,fill=protowhite,inner sep=1.2pt] {frozen at} (m4);
\draw[gvedge] (r9.east) to[bend left=18]
  node[font=\tiny,fill=protowhite,inner sep=1.2pt,pos=0.55] {frozen at} (m4.south west);
\draw[gvedge] (r6) -- node[font=\tiny,fill=protowhite,inner sep=1.2pt] {verified at} (m3);
\draw[gvedge] (r7) -- node[font=\tiny,fill=protowhite,inner sep=1.2pt] {cleared at} (m8);
\node[font=\tiny\itshape,text=protogray,anchor=north west,text width=132mm]
  at (0,-7.35) {Every edge carries its link type, because naming the type is
  what stops the chain being read as an inheritance. A tier 1 result does not
  become a tier 3 result by travelling along an arrow.};
\end{appfig}
\vspace{-0.65cm}
\figcaption{Five published quantities and the exact filing section each one\\
reaches. Every edge carries its link type, and every chain ends\\
at a milestone whose artifact a program officer can request.}
\end{appfloat}

Five link types appear on the fifteen edges, and each licenses something
different. \emph{Cited} means a published result is quoted with its own interval
and limitation, and does not license restating a metastatic result as a
resectable one. \emph{Informs} means a simulated result shaped a design choice,
and does not license presenting a simulated survival figure as an expected
outcome. \emph{Supports} means a verified computation raises confidence in a
stated assumption, and does not substitute for a clinical measurement.
\emph{Specifies} means a bench measurement fixes a numeric requirement, and does
not claim the requirement has been met in a patient. \emph{Bounds} means a
result sets a monitoring limit, and does not make that limit a predicted rate.

\subsection{What the Work Cost, and What It Costs Elsewhere}

\begin{apptable}
\begin{tabularx}{\textwidth}{>{\raggedright\arraybackslash}X >{\raggedright\arraybackslash}p{2.9cm} >{\raggedright\arraybackslash}p{3.2cm} >{\raggedright\arraybackslash}p{3.0cm}}
\headrow \thead{Work product} & \thead{Industry cost} & \thead{Industry schedule} & \thead{This author}\\
\midrule
Empirical virtual trial, triplicate & above \$120,000 per run & 4.5 months & \$36,330 estimated, about one month \\
Ten-arm QSP virtual trial & above \$2,000,000 & several months to over a year & \$36,330 estimated, about one month \\
Digital-twin simulator platform & \$28,000 to \$700,000 & 3 to 16 weeks & \$36,330 estimated, about one month \\
\end{tabularx}
\end{apptable}
\vspace{-0.65cm}
\tabcap{Three work products against their industry cost and schedule benchmarks.
The author figures are estimates from the source set and are reported as such,
not as audited accounts.}

The same cost discipline that produced a \$36,330 virtual trial produces a 7.5
percent indirect rate. Both are the same fact seen twice: a firm with no
facilities, no administrative layer and no negotiated rate agreement has very
little between an appropriation and the work it buys. Figure 3 prices that fact
at \$531,000 against a university route.
